\talkslide{alternation}{The cat-and-mouse game problem}{Timing: provisional}{The conventional loop alternates substantial training phases.}{
\advprompterflow{2}
}
\speak{The conventional loop alternates substantial training phases. Each side improves against a snapshot of its opponent. This can be useful, but it may spend effort on outdated attacks or an outdated target, and it does not by itself guarantee convergence. The problem we want to address is the mismatch between a moving model and a relatively static training challenge. AdvPrompter provides the concrete example of the first two steps: learn an attacker against a fixed target, then use its generated attacks for defense fine-tuning. For a static attack dataset used in RL, refer back to the Beam safety recipe: sample a response to a collected attack, judge it in context, and update the defender. The Iterate step describes the general cat-and-mouse pattern, not a claim that AdvPrompter jointly updates both models throughout training.}
